% !TEX root = ../main.tex

\chapter{Introduction to Web Design and HTML}

\section{Core Concepts in Web Design}

\begin{definitionbox}[Internet]
    \term{Internet} is the connection of millions of computers worldwide that allows them to communicate and share information.
\end{definitionbox}

\begin{definitionbox}[World Wide Web]
    \term{World Wide Web} is a system of sharing information over the Internet by using hypertext documents, which are linked together and can be accessed through a web browser.
\end{definitionbox}

\Note The Internet is the hardware, wires and protocols: the global infrastructure. The World Wide Web is a service, like email or online gaming, that runs on top of that infrastructure.

\subsection{Basic Components of the Web}

The web works through the combined use of three main components:

\begin{enumerate}
    \item \term{HTML}: A markup language used to write and structure webpages.
    \item \term{HTTP}: The protocol used to transfer webpages between web server and browser.
    \item \term{Web Browser}: The software that receives, interprets and displays webpage content for the user.
\end{enumerate}

\par\vspace{\baselineskip}

\begin{definitionbox}[Hypertext document file]
    A hypertext document is a type of document that contains links to other documents, allowing users to navigate between them easily. 
    \begin{itemize}
        \item Like images this is a file type with extension .html or .htm.
        \item It is made by using HTML (HyperText Markup Language).
    \end{itemize}
\end{definitionbox}

\begin{definitionbox}[Webpage]
    A webpage is a single hypertext document or page that is displayed in a web browser. It is also called a web document. It can contain text, audio, video, still images, animation, links of other webpages and other multimedia content.
\end{definitionbox}

\begin{definitionbox}[Website]
    A website is a collection of related webpages stored under a common domain name and connected through hyperlinks.
\end{definitionbox}

\begin{definitionbox}[Web design]
    Web design is the process of planning, creating, arranging and presenting information on the World Wide Web. It includes the design of webpages, navigation system, content layout, links, images, and user interaction.
\end{definitionbox}

\subsection{Basic Relationship}

\begin{itemize}
    \item A webpage is one single page.
    \item A website contains multiple related webpages.
    \item The first page shown after entering a website is called the \term{homepage}.
    \item Webpages are connected using hyperlinks.
    \item A web browser is used to view webpages.
\end{itemize}

\begin{summarybox}[Core Idea]
    HTML $\rightarrow$ Hypertext document files $\rightarrow$ Webpage $\xrightarrow{\text{Web Design}}$ Website $\rightarrow$ World Wide Web $\rightarrow$ Internet
\end{summarybox}


\subsection{Advantages of Website}

\begin{itemize}
    \item Information can be published and accessed easily.
    \item A website can be accessed from anywhere in the world.
    \item It is available twenty-four hours a day.
    \item It reduces the need for physical communication.
    \item It supports text, images, audio, video and links.
    \item Information can be updated more easily than printed documents.
    \item It is useful for education, business, news, communication and entertainment.
\end{itemize}


\section{Website}

\subsection{Website Structure}

\begin{definitionbox}[Website Structure]
Website structure refers to the way webpages are organized and interconnected within a website.
\end{definitionbox}

\subsubsection{Common Pages in a Website}

\begin{itemize}
    \item \term{Home Page}: the first page displayed after entering a website. It introduces the website and usually provides links to main sections.
    \item \term{Main Section Page}: a page that represents a major topic or section of the website.
    \item \term{Subsection Page}: a page placed under a main section to present more specific information.
\end{itemize}

\subsubsection{Types of Website Structure}

\paragraph{Hierarchical Structure}

\begin{itemize}
    \item Pages are arranged like a tree.
    \item Starts from a homepage and branches into sections.
    \item Each section may contain sub-sections.
    \item Easy to manage and most commonly used.
\end{itemize}

\paragraph{Linear Structure}

\begin{itemize}
    \item Pages are connected in a sequence.
    \item Users move step-by-step (next or previous).
    \item Suitable for tutorials or guided content.
\end{itemize}

\paragraph{Network Structure}

\begin{itemize}
    \item Each page is connected with many other pages.
    \item No fixed navigation path.
    \item Provides flexibility but may confuse users.
\end{itemize}

\paragraph{Hybrid Structure}

\begin{itemize}
    \item Combination of multiple structures.
    \item Uses hierarchical with linear or network links.
    \item Most modern websites follow this model.
\end{itemize}

\begin{summarybox}
A good website structure improves navigation, usability, and overall performance.
\end{summarybox}


% ─────────────────────────────────────────
\subsection{Types of Website}
% ─────────────────────────────────────────

\subsubsection{Static Website}

\begin{definitionbox}[Static Website]
A static website displays fixed content that does not change unless manually updated.
\end{definitionbox}

\begin{itemize}
    \item Built using HTML and CSS.
    \item Content remains the same for all users until the source files are manually edited.
    \item Users cannot directly input or update information on the page.
    \item No database connection is normally required.
    \item Communication is mainly one-way: server to client.
    \item Faster loading speed.
    \item Simple to develop.
\end{itemize}

\paragraph{Advantages of Static Website}

\begin{itemize}
    \item Low development and hosting cost.
    \item Easy to develop, test and control.
    \item Loads quickly because pages are already prepared.
    \item More secure because there is usually no database or server-side processing.
\end{itemize}

\paragraph{Disadvantages of Static Website}

\begin{itemize}
    \item Content cannot change automatically.
    \item User input or update is not supported.
    \item Updating many pages manually can take time.
    \item Not suitable for large interactive services.
\end{itemize}

\subsubsection{Dynamic Website}

\begin{definitionbox}[Dynamic Website]
A dynamic website generates content dynamically based on user interaction or database.
\end{definitionbox}

\begin{itemize}
    \item Uses server-side scripting such as PHP, ASP.NET, JSP or similar technologies.
    \item Content can change for different users.
    \item Runtime page content or design can change based on request, time, user or database.
    \item Supports database connection.
    \item Users can provide input and update information.
    \item More complex than static websites.
\end{itemize}

\paragraph{Advantages of Dynamic Website}

\begin{itemize}
    \item Supports interactive services such as login, search, comment and form submission.
    \item Can display different content for different users.
    \item Handles large amounts of information with database support.
    \item Easier to update common content from a central system.
\end{itemize}

\paragraph{Disadvantages of Dynamic Website}

\begin{itemize}
    \item Development and maintenance are more complex.
    \item Cost is higher than static websites.
    \item Loading may be slower because pages are generated at runtime.
    \item Security management is more important because input and database are involved.
\end{itemize}

\subsubsection{Comparison}

\begin{longtable}{|C{3cm}|C{6.5cm}|C{6.5cm}|}
\longtableheader{3}{Static vs Dynamic Website}{%
\tblcolhdcell{Feature} & \tblcolhdcell{Static Website} & \tblcolhdcell{Dynamic Website}}

Content & Fixed & Changeable at runtime \\ \tblhline
Update & Manual file editing & Automatic or database-driven update \\ \tblhline
User input & Not supported & Supported \\ \tblhline
Database & Not required & Usually required \\ \tblhline
Communication & Mainly server to client & Client-server interaction \\ \tblhline
Speed & Faster & Slower than static pages \\ \tblhline
Technology & HTML, CSS & PHP, ASP.NET, JSP, database \\ \tblhline
Complexity & Simple & Complex \\ \tblhline
Cost & Lower & Higher \\ \tblhline

\tblbottomrule
\caption{Comparison between Static and Dynamic Website}
\end{longtable}

\subsection{Types of Webpage Based on Location}

\begin{itemize}
    \item \term{Local Webpage}: a webpage stored and opened from the user's own computer or local storage. It can be viewed without uploading it to a web server.
    \item \term{Remote Webpage}: a webpage stored on a remote web server and accessed through the internet using a URL.
\end{itemize}

\subsection{Website Classification by Purpose}

\begin{itemize}
    \item \term{Archive Site}: Stores old or important information for later use.
    \item \term{Business Site}: Presents information about a business, product or service.
    \item \term{E-commerce Site}: Supports buying and selling products or services online.
    \item \term{Community Site}: Connects people with common interests.
    \item \term{Database Site}: Provides searchable or organized data from a database.
    \item \term{Development Site}: Shares software, tools, documentation or project resources.
    \item \term{Directory Site}: Lists links, organizations, services or resources by category.
    \item \term{Download Site}: Provides files, software, documents or media for download.
    \item \term{Game Site}: Provides online games or gaming-related content.
    \item \term{Information Site}: Publishes subject-based information for readers.
    \item \term{News Site}: Publishes current news and updates.
\end{itemize}

\begin{summarybox}
Static websites are simple and fast, whereas dynamic websites are flexible and interactive.
\end{summarybox}

% ─────────────────────────────────────────
\section{Web Browser and Search Engine}
% ─────────────────────────────────────────

\subsection{Web Browser}

\begin{definitionbox}[Web Browser]
A web browser is application software used to access, retrieve, interpret and display webpages.
\end{definitionbox}

A browser works like an interpreter. It receives webpage code from the web server, converts it into visual output and shows the webpage on the screen.

\subsubsection{Functions of Web Browser}

\begin{itemize}
    \item Opens websites and webpages using a URL.
    \item Interprets HTML and other webpage code.
    \item Displays text, images, links, forms and other webpage contents.
    \item Helps users move from one webpage to another through hyperlinks.
    \item Can download information or files from websites.
    \item Can upload information through forms or file upload systems.
    \item Can use search engines when the exact web address is not known.
    \item Example: Google Chrome, Mozilla Firefox, Microsoft Edge, Safari, Opera.
\end{itemize}

\subsection{Web Browsing}

\begin{definitionbox}[Web Browsing]
Web browsing is the process of visiting and viewing webpages stored on different web servers through a web browser.
\end{definitionbox}

\subsubsection{Methods of Web Browsing}

\begin{itemize}
    \item \term{Direct Browsing by URL}: The user types the known website address in the browser address bar and opens the webpage directly.
    \item \term{Browsing through Search Engine}: The user searches with keywords when the exact website address is not known.
\end{itemize}


\subsection{Search Engine}

\begin{definitionbox}[Search Engine]
A search engine is a tool that helps users find information on the internet by searching indexed webpages and showing relevant results.
\end{definitionbox}

A search engine collects and organizes webpage information. When a user searches with a word or phrase, it shows a list of related webpages.

\subsubsection{Functions of Search Engine}

\begin{itemize}
    \item Helps users find information when the exact website address is unknown.
    \item Searches indexed webpages using keywords.
    \item Displays related webpage links as search results.
    \item Helps users reach websites quickly through the result list.
    \item Example: Google, Bing, Yahoo, DuckDuckGo, Ask, Yandex, Baidu, WolframAlpha (computational information-based search tool).
\end{itemize}

\subsection{Difference Between Browser and Search Engine}

\begin{longtable}{|C{2.5cm}|L{6.5cm}|L{6.5cm}|}
\longtableheader{3}{Browser vs Search Engine}{%
\tblcolhdcell{Aspect} & \tblcolhdcell{Web Browser} & \tblcolhdcell{Search Engine}}

Meaning
& Software used to open and display webpages.
& Tool used to find information and webpages on the internet. \\ \tblhline

Main Work
& Reads webpage code and shows the webpage visually.
& Searches indexed webpages and shows related results. \\ \tblhline

Use
& Used when opening a website or viewing a webpage.
& Used when the needed website or exact address is not known. \\ \tblhline

Input
& URL or website address.
& Keyword, phrase or question. \\ \tblhline

Examples
& Google Chrome, Mozilla Firefox, Microsoft Edge, Safari, Opera.
& Google, Bing, Yahoo, DuckDuckGo, Ask. \\ \tblhline

\tblbottomrule
\caption{Difference between Web Browser and Search Engine}
\label{tab:browser-search-engine}
\end{longtable}

\begin{summarybox}
A web browser opens and displays webpages, while a search engine helps find webpages and information.
\end{summarybox}

\section{IP Address}

\begin{definitionbox}[IP Address]
IP Address stands for Internet Protocol Address. It is a unique numerical address used to identify a device connected to a network on the internet.
\end{definitionbox}

\Note An IP address works like the digital address of a device. When sender sends data to a receiver, the receiver's IP address is used to ensure that the data reaches the correct destination.

\subsection{Purpose of IP Address}

\begin{itemize}
    \item It identifies a device in a network.
    \item It helps send data to the correct destination.
    \item It allows communication between client and server.
    \item It is necessary for internet-based communication.
\end{itemize}

\subsection{Types of IP Address}

There are two main types of IP addresses:
\begin{itemize}
    \item \term{IPv4}: 32-bit address, with four 8-bit octets, 1 octet = 8 bits, usually written in four decimal parts.\\
          Example:  \texttt{11000000.10101000.00000001.00000001} $\xrightarrow{\text{in decimal}}$ \texttt{192.168.1.1}
    \item \term{IPv6}: 128-bit address, with eight 16-bit words, usually written in hexadecimal format.\\
          Example: \texttt{2001:0db8:85a3:0000:0000:8a2e:0370:7334}
\end{itemize}

\subsubsection{Network ID and Host ID}
\begin{itemize}
    \item The internet is a network of many networks.
    \item Each network has many hosts (computers, mobiles, servers etc.) in it.
    \item So each network and each device needs an unique identifier for communication.
    \item This is done by dividing the IP address into two parts: Network ID and Host ID.
        \begin{itemize}
            \item \term{Network ID}: identifies the network the device is connected to. \\
            Example: in \texttt{192.168.1.1}, the network ID is \texttt{192.168.1}.
            \item \term{Host ID}: identifies the specific device inside that network. \\
            Example: in \texttt{192.168.1.1}, the host ID is \texttt{1}.
        \end{itemize}
\end{itemize}

\begin{questionanswerbox}
\textbf{Q: Why is IP address required?}

\medskip
\textbf{A:}  
IP address is required because every device in a network must be uniquely identified. It helps data packets travel from the sender to the correct receiver.

\medskip
\textbf{Conclusion:}  
Without IP addresses, devices on the internet could not communicate properly.
\end{questionanswerbox}


\section{Domain Name and Domain Name System (DNS)}

\begin{definitionbox}[Domain Name]
    A \term{domain name} is a human-readable name used to identify a website on the internet.
\end{definitionbox}

\Note IP addresses are difficult for humans to remember. For this reason, domain names act as an alias for the numeric IP address.

\begin{definitionbox}[DNS]
    \term{DNS} (Domain Name System) is the system that converts a readable domain name into its corresponding numeric IP address so the browser can load the resource.
\end{definitionbox}

\begin{summarybox}[The Bridge]
    \begin{itemize}
        \item \term{IP Address} ($192.168.1.1$): used by computers in numeric form.
        \item \term{Domain Name} ($google.com$): used by humans in readable form.
        \item \term{DNS}: works like a phonebook that connects the readable name with the numeric address.
    \end{itemize}
\end{summarybox}

\subsection{Some Technical Terms}
\begin{itemize}
    \item \term{FQDN} (Fully Qualified Domain Name): The complete domain name that identifies a specific host on the internet.
    \item \term{Subdomain}: An additional name placed before the main domain name, such as \texttt{www}.
    \item \term{Web Hosting}: The server space where website files are stored so users can access them through the internet.
    \item \term{Difference}: The \textit{Domain} is your address; \textit{Hosting} is the house where your data lives.
\end{itemize}

\subsection{Domain Registration and ICANN}
\begin{itemize}
    \item \term{Domain Registration} is the process of reserving a domain name for use on the internet. It involves choosing a unique name and registering it with a domain registrar.
    \item Every registered domain name must be unique. If the same name could be registered by many people, users would not know which website to visit. 
    \item For this reason, domain names are controlled through registration authorities.
    \item \term{ICANN} stands for "Internet Corporation for Assigned Names and Numbers".
    \item ICANN coordinates the global domain name and IP address system.
    \item Domain names are usually registered through domain registrars approved under this system.
    \item A registered domain gives a website a unique identity on the internet.
\end{itemize}

\subsection{Types of Domain}

\begin{itemize}
    \item \textbf{Top Level Domain (TLD)} $\rightarrow$ The last part of a domain name. It shows the highest level/category of the domain.
    \begin{itemize}
        \item \textbf{gTLD (Generic TLD)} $\rightarrow$ General TLD, not tied to a country. Examples: \texttt{.com}, \texttt{.org}, \texttt{.net}
        \item \textbf{ccTLD (Country Code TLD)} $\rightarrow$ Country-based TLD. Examples: \texttt{.bd}, \texttt{.uk}, \texttt{.in}
    \end{itemize}

    \item \textbf{Second Level Domain (SLD) / Main Domain Name} $\rightarrow$ The unique name chosen by the owner. It identifies the website or organization. Examples: \texttt{google}, \texttt{bbc}

    \item \textbf{Subdomain / Host Name} $\rightarrow$ A name placed before the main domain to identify a service or section. Examples: \texttt{www}, \texttt{mail}, \texttt{blog}
\end{itemize}

\begin{infobox}[Example: Domain Structure]

\[
\small
\underbrace{\text{\texttt{mail}}}_{\substack{\text{Subdomain}\\\text{/ Host Name}}}
\text{\texttt{.}}
\underbrace{\text{\texttt{google}}}_{\substack{\text{Main Domain}\\\text{/ SLD}}}
\underbrace{\text{\texttt{.com}}}_{\substack{\text{TLD}\\\text{gTLD}}}
\]

\[
\small
\underbrace{\text{\texttt{www}}}_{\substack{\text{Subdomain}\\\text{/ Host Name}}}
\text{\texttt{.}}
\underbrace{\text{\texttt{bbc}}}_{\substack{\text{Main Domain}\\\text{/ SLD}}}
\underbrace{\text{\texttt{.co}}}_{\substack{\text{Commercial}\\\text{category}}}
\underbrace{\text{\texttt{.uk}}}_{\substack{\text{TLD}\\\text{ccTLD}}}
\]

\end{infobox}

\subsection{Common Top Level Domains / Generic TLDs}

\begin{longtable}{|C{3cm}|C{6cm}|}
\longtableheader{2}{Common Top Level Domains}{%
\tblcolhdcell{Domain} & \tblcolhdcell{Used For}}

.com & Commercial organizations \\ \tblhline
.net & Network-related services \\ \tblhline
.edu & Educational institutions \\ \tblhline
.gov & Government organizations \\ \tblhline
.org & Non-profit organizations \\ \tblhline
.mil & Military organizations \\ \tblhline
.int & International organizations \\ \tblhline
.biz & Business organizations \\ \tblhline
.info & Information-based websites \\ \tblhline

\tblbottomrule
\caption{Common Top Level Domains}
\label{tab:common-tld}
\end{longtable}

\subsection{Country Code Top Level Domains}

\begin{longtable}{|C{3cm}|C{4cm}|C{3cm}|C{4cm}|}
\longtableheader{4}{Common Country Code Top Level Domains}{%
\tblcolhdcell{Code} & \tblcolhdcell{Country} & \tblcolhdcell{Code} & \tblcolhdcell{Country}}

.bd & Bangladesh & .cn & China \\ \tblhline
.au & Australia & .fi & Finland \\ \tblhline
.in & India & .fr & France \\ \tblhline
.br & Brazil & .hk & Hong Kong \\ \tblhline
.uk & United Kingdom & .us & United States \\ \tblhline

\tblbottomrule
\caption{Common Country Code Top Level Domains}
\label{tab:common-cctld}
\end{longtable}

\section{Web Hosting}

\begin{definitionbox}[Web Hosting]
Web hosting means storing website files on an internet-connected server so that the website can be uploaded and accessed online.
\end{definitionbox}

\begin{itemize}
    \item Hosting space is usually provided by an ISP or hosting company.
    \item Hosting cost depends on storage space, bandwidth, server resources and support.
    \item Different hosting services can be obtained from hosting providers by paying rent or service charge.
    \item Without hosting, a website's files cannot be publicly accessed through the internet.
\end{itemize}

\subsection{Types of Web Hosting}

\subsubsection{Based on Platform or Operating System}

\begin{itemize}
    \item \term{Windows Hosting}: Hosting on a Windows server. It is mainly used for websites made with ASP/ASP.NET and Microsoft SQL Server.
    
    \item \term{Linux Hosting}: Hosting on a Linux server. It is mainly used for websites made with PHP and MySQL.
\end{itemize}

\subsubsection{Based on Facilities}

\begin{itemize}
    \item \term{Free Hosting}: Basic hosting for small personal websites. It costs little or nothing, but bandwidth, security and domain facilities are limited.
    
    \item \term{Shared Hosting}: Many websites share the same server resources. It is cheaper, but speed, security, email, database and bandwidth are limited.
    
    \item \term{Dedicated Hosting}: One website or user gets a full server. It gives better speed, security, bandwidth, database support and control, so it is suitable for large websites.
    
    \begin{itemize}
        \item \term{Managed Hosting}: The hosting company manages setup, software, security, maintenance and control panel support. It is easier but more costly.
        \item \term{Unmanaged Hosting}: The website owner manages setup, software, security and maintenance. It is cheaper but needs technical skill.
    \end{itemize}
    
    \item \term{Collocated Hosting}: The owner keeps their own server in the hosting company's data center and uses the company's power, network and support.
\end{itemize}

\subsubsection{Email Services with Hosting}

Hosting providers may also provide domain-based email accounts, such as \texttt{info@example.com}.

\begin{itemize}
    \item \term{POP} (Post Office Protocol): downloads email from the server to a client device.
    \item \term{IMAP} (Internet Message Access Protocol): keeps email on the server and synchronizes it across devices.
    \item \term{Web-based Email / Webmail}: allows users to read and send email through a web browser.
\end{itemize}

\subsection{Difference between Domain Name and Web Hosting}

\begin{longtable}{|L{8.5cm}|L{8.5cm}|}
\longtableheader{2}{Domain Name vs Web Hosting}{%
\tblcolhdcell{Domain Name} & \tblcolhdcell{Web Hosting}}

The unique readable name of a website. & The server space where website files are stored. \\ \tblhline
Must be registered with a domain registrar. & Must be rented or configured from a hosting provider. \\ \tblhline
One domain name identifies one main web address. & One hosting account may contain one or many websites, depending on the plan. \\ \tblhline
Users cannot easily reach a website without a domain name or URL. & A website cannot be publicly available without hosting. \\ \tblhline
Domain registration does not require server configuration. & Hosting requires server configuration, storage, bandwidth and control panel settings. \\ \tblhline
Bandwidth is not the main issue in domain registration. & Bandwidth is an important issue in hosting. \\ \tblhline
The quality of the domain registrar is less visible to normal visitors. & The quality of the hosting provider affects speed, uptime, security and support. \\ \tblhline

\tblbottomrule
\caption{Difference between Domain Name and Web Hosting}
\label{tab:domain-hosting-difference}
\end{longtable}

% ─────────────────────────────────────────
\section{URL (Uniform Resource Locator)}
% ─────────────────────────────────────────

\begin{definitionbox}[URL]
A URL (Uniform Resource Locator) is the complete web address used to locate a resource on the internet.
\end{definitionbox}

\Note A URL identifies where a webpage or file is located and also tells the browser which protocol should be used to access it.

\subsection{Components of URL}

\begin{itemize}
    \item \term{Protocol}: Defines the communication rule used to access the resource. Examples: \texttt{http://}, \texttt{https://}, \texttt{ftp://}
    \item \term{Domain Name}: Identifies the website or server. Example: \texttt{www.example.com}
    \item \term{Path}: Specifies the exact location of a file or page inside the server. Example: \texttt{/folder/page.html}
\end{itemize}

\begin{formulabox}[Structure of URL]
\[
\text{URL} = \text{Protocol} + \text{Domain Name} + \text{Path}
\]
\[
\text{URL}
=
\underbrace{\text{\texttt{https://}}}_{\text{Protocol}}
\underbrace{\text{\texttt{www.example.com}}}_{\text{Domain Name}}
\underbrace{\text{\texttt{/folder/page.html}}}_{\text{Path}}
\]
\end{formulabox}

\begin{infobox}[Example: URL Structure]

\[
\small
\underbrace{\text{\texttt{https://}}}_{\text{Protocol}}
\underbrace{\text{\texttt{www}}}_{\substack{\text{Subdomain}\\\text{/ Host Name}}}
\text{\texttt{.}}
\underbrace{\text{\texttt{example}}}_{\substack{\text{Main Domain}\\\text{/ SLD}}}
\underbrace{\text{\texttt{.com}}}_{\text{TLD}}
\underbrace{\text{\texttt{/folder/page.html}}}_{\text{Path}}
\]

\[
\small
\underbrace{\text{\texttt{http://}}}_{\text{Protocol}}
\underbrace{\text{\texttt{www}}}_{\substack{\text{Subdomain}\\\text{/ Host Name}}}
\text{\texttt{.}}
\underbrace{\text{\texttt{corneliusenterprises}}}_{\substack{\text{Main Domain}\\\text{/ SLD}}}
\underbrace{\text{\texttt{.co}}}_{\substack{\text{Commercial}\\\text{category}}}
\underbrace{\text{\texttt{.uk}}}_{\substack{\text{Country Code}\\\text{TLD}}}
\]

\end{infobox}

\begin{summarybox}
A URL gives both the access method and the location of a web resource.
\end{summarybox}

% ─────────────────────────────────────────
\section{Protocol}
% ─────────────────────────────────────────

\begin{definitionbox}[Protocol]
A protocol is a set of rules that controls how data is transmitted between computers on the internet.
\end{definitionbox}

\subsection{Common Internet Protocols}

\begin{itemize}
    \item \term{HTTP} (Hyper Text Transfer Protocol): Used to transfer webpages between browser and web server.
    \item \term{FTP} (File Transfer Protocol): Used to upload and download files through the internet.
    \item \term{IP} (Internet Protocol): Handles addressing and routing so data packets reach the correct destination.
    \item \term{TCP/IP} (Transmission Control Protocol / Internet Protocol): Main protocol system used for internet communication.
    \item \term{SMTP} (Simple Mail Transfer Protocol): Used for sending email.
    \item \term{POP} (Post Office Protocol): Used for receiving or downloading email from a mail server.
    \item \term{VoIP} (Voice over Internet Protocol): Used for voice communication through the internet.
\end{itemize}

\subsection{HTTP}

\begin{definitionbox}[HTTP]
HTTP stands for Hyper Text Transfer Protocol. It is the web protocol used to exchange data between a browser and a web server.
\end{definitionbox}

\Note HTTP is used when a user visits or browses a webpage. The browser sends a request to the web server, and the server sends the requested webpage or resource back to the browser.

\subsubsection{Functions of HTTP}

\begin{itemize}
    \item Creates a connection between the browser and the server.
    \item Sends the browser's request to the server.
    \item Brings the requested webpage or resource from the server to the browser.
    \item Disconnects the browser and server after the request-response process.
\end{itemize}

\subsection{FTP}

\begin{definitionbox}[FTP]
FTP stands for File Transfer Protocol. It is a protocol used to exchange files between computers through the internet.
\end{definitionbox}

\Note FTP is mainly used to upload and download files between a local computer and a remote computer or server. It is commonly used to upload website files to a hosting server.

\subsubsection{Functions of FTP}

\begin{itemize}
    \item Transfers website files, documents, books, images, software and other normal files.
    \item \term{Download}: Receives a file from a remote computer or server to the local computer.
    \item \term{Upload}: Sends a file from the local computer to a remote computer or server.
    \item FTP software such as \term{WS-FTP} and \term{Fetch} can be used for upload and download.
\end{itemize}

\subsubsection{Types of FTP Access}

\begin{itemize}
    \item \term{Private Access}: Requires username and password to access server files.
    \item \term{Public Access}: Allows users to access selected files without private login.
\end{itemize}

\begin{summarybox}[HTTP vs FTP]
HTTP is used to transfer hypertext/webpage files between browser and web server. FTP is used to transfer normal files such as documents, images, software and website files between computers.
\end{summarybox}

% ─────────────────────────────────────────
\section{Web Communication}
% ─────────────────────────────────────────

\begin{definitionbox}[Web Communication]
Web communication is the process of exchanging data between a client browser and a web server through the internet.
\end{definitionbox}

\Note The World Wide Web works on top of the internet using a client-server system. The browser works as the client, and the web server stores and sends the requested webpage or resource.

\subsection{Correct Overall Idea}

\begin{itemize}
    \item \term{Internet}: The global network infrastructure that connects computers and servers.
    \item \term{World Wide Web}: A service that runs on the internet and uses webpages, hyperlinks, browsers and web servers.
    \item \term{Website Developer}: Creates website files using HTML, CSS, images, scripts and other resources.
    \item \term{Web Hosting}: Stores the website files on an internet-connected server.
    \item \term{IP Address}: The numerical address of the hosting server.
    \item \term{Domain Name}: The human-readable name registered through a domain registrar.
    \item \term{DNS}: Connects the domain name with the hosting server's IP address.
    \item \term{URL}: The complete address used by the browser to request a specific webpage or file.
    \item \term{Protocol}: The rule used for communication, such as HTTP or HTTPS.
\end{itemize}

\subsection{Website Creation Side}

\begin{enumerate}
    \item The developer creates hypertext documents using \term{HTML}.
    \item Other files such as CSS, images, audio, video and scripts may also be added.
    \item These files are organized together as a website project or website folder.
    \item The website folder is uploaded to a web hosting server.
    \item The hosting server has an IP address.
    \item A domain name is registered through a domain registrar.
    \item DNS records connect the domain name with the hosting server's IP address.
\end{enumerate}

\subsection{User Access Side}

\begin{enumerate}
    \item The user enters a URL in the web browser.
    \item The browser reads the protocol, domain name and path from the URL.
    \item DNS converts the domain name into the server's IP address.
    \item The browser connects to the correct web server using HTTP or HTTPS.
    \item The browser sends an HTTP request for the specific webpage or file.
    \item The server searches for the requested HTML document or resource.
    \item The server sends the requested data back as an HTTP response.
    \item The browser interprets the HTML and related files.
    \item The webpage is displayed on the screen.
\end{enumerate}

\begin{summarybox}[Web Working Flow]
Developer creates website files $\rightarrow$ files are uploaded to hosting server $\rightarrow$ domain name is connected to server IP through DNS $\rightarrow$ user enters URL $\rightarrow$ browser requests the file $\rightarrow$ server responds $\rightarrow$ browser displays the webpage.
\end{summarybox}

\section{Difference between IP Address, Domain Name, DNS and URL}
\begin{longtable}{|C{2cm}|L{3.4cm}|L{3.4cm}|L{3.4cm}|L{3.4cm}|}
\longtableheader{5}{IP Address vs Domain Name vs DNS vs URL}{%
\tblcolhdcell{Aspect} & \tblcolhdcell{IP Address} & \tblcolhdcell{Domain Name} & \tblcolhdcell{DNS} & \tblcolhdcell{URL}}

Meaning
& A unique numeric address used to identify a device or server on a network.
& A human-readable name used to identify a website.
& A system that converts a domain name into its corresponding IP address.
& A complete web address used to locate a specific resource on the internet. \\ \tblhline

Main User
& Computers and network devices.
& Humans.
& Browser and network system.
& Browser and user. \\ \tblhline

Main Work
& Finds the actual destination server or device.
& Gives an easy readable name for a website.
& Connects the domain name with the correct IP address.
& Gives the access method and exact resource location. \\ \tblhline

Points To
& A device or server.
& A website name.
& The IP address of a domain.
& A webpage, file, image or resource. \\ \tblhline

Form
& Numeric form.
& Readable word form.
& Background lookup system.
& Full address with protocol, domain and path. \\ \tblhline

Example
& \texttt{142.250.190.78}
& \texttt{google.com}
& \texttt{google.com} $\rightarrow$ \texttt{142.250.190.78}
& {\hypersetup{urlcolor=black}\footnotesize\url{https://www.example.com/folder/page.html}} \\ \tblhline

\tblbottomrule
\caption{Difference between IP Address, Domain Name, DNS and URL}
\label{tab:ip-domain-dns-url}
\end{longtable}

% ─────────────────────────────────────────
\section{HTML (HyperText Markup Language)}
% ─────────────────────────────────────────

\begin{definitionbox}[HTML]
HTML (HyperText Markup Language) is the standard markup language used to create and structure webpages.
\end{definitionbox}

HTML does not perform computation; it describes the structure and content of a webpage so that a browser can render it.

\subsection{Short History of HTML}

\begin{itemize}
    \item Tim Berners-Lee introduced HTML while working at CERN in 1990.
    \item HTML 2.0 was released in 1995.
    \item HTML 3.2 was released in January 1997.
    \item HTML 4 was released in December 1997.
    \item HTML5 is the modern version used for current web development.
\end{itemize}

\subsection{Key Characteristics}

\begin{itemize}
    \item Uses tags to mark content.
    \item Works with web browsers to display pages.
    \item Supports text, images, links, lists, tables and forms.
    \item Platform independent (runs in any modern browser).
\end{itemize}

\subsection{Advantages of HTML}

\begin{itemize}
    \item Simple and easy to learn.
    \item Supported by all browsers.
    \item Integrates with CSS and JavaScript.
    \item Lightweight and fast to load.
    \item Can be written in any simple text editor.
    \item Does not require a separate license to use.
\end{itemize}

\subsection{Limitations of HTML}

\begin{itemize}
    \item Cannot create dynamic behavior by itself.
    \item Limited styling without CSS.
    \item No database interaction.
    \item Large pages may require many lines of code.
    \item Security and data processing must be handled by other technologies.
\end{itemize}

\begin{summarybox}
HTML defines structure; CSS handles styling; JavaScript adds interactivity.
\end{summarybox}


% ─────────────────────────────────────────
\section{Basic Structure of an HTML Document}
% ─────────────────────────────────────────

\begin{definitionbox}[HTML Document Structure]
An HTML document is organized using a standard structure consisting of \texttt{<!DOCTYPE>}, \texttt{<html>}, \texttt{<head>} and \texttt{<body>}.
\end{definitionbox}

\subsection{Skeleton Code}

This example shows the minimum structure of a valid HTML page.
The browser uses \texttt{<!DOCTYPE html>} to understand that the document follows HTML5 rules.
The \texttt{head} stores page information, while the \texttt{body} stores visible content.

\begin{lstlisting}[style=htmlide, caption={Basic HTML Document Structure}]
<!DOCTYPE html>
<html>
<head>
    <title>Page Title</title>
</head>
<body>
    Content goes here
</body>
</html>
\end{lstlisting}

\textbf{What to notice}
\begin{itemize}
    \item The \texttt{html} element contains the whole page.
    \item The \texttt{title} appears in the browser tab.
    \item Visible page content is written inside \texttt{body}.
\end{itemize}

\subsection{Explanation of Parts}

\begin{itemize}
    \item \texttt{<!DOCTYPE html>} declares the document type (HTML5).
    \item \texttt{<html>} is the root element.
    \item \texttt{<head>} contains metadata (title, links, scripts).
    \item \texttt{<title>} sets the browser tab title.
    \item \texttt{<body>} contains visible content.
\end{itemize}

\subsection{Head Section Elements}

The head section contains information about the document. Most of this information is not directly displayed as page content.

\begin{longtable}{|C{3cm}|C{3cm}|L{8cm}|}
\longtableheader{3}{Common Head Section Elements}{%
\tblcolhdcell{Element} & \tblcolhdcell{Type} & \tblcolhdcell{Use}}

\texttt{title} & Container & Sets the page title shown in the browser tab \\ \tblhline
\texttt{meta} & Empty & Stores document information such as author, charset or keywords \\ \tblhline
\texttt{link} & Empty & Connects external files such as CSS stylesheets \\ \tblhline
\texttt{base} & Empty & Sets a base URL for relative links \\ \tblhline
\texttt{style} & Container & Stores internal CSS rules \\ \tblhline
\texttt{script} & Container & Stores or connects script commands \\ \tblhline

\tblbottomrule
\caption{Common elements used in the HTML head section}
\end{longtable}

\subsection{Body Section Elements}

The body section begins after the closing \texttt{</head>} tag. It contains the visible text, images, links, lists, tables and other contents of a webpage.

\begin{itemize}
    \item Block-level elements organize larger parts of a page.
    \item Text-level elements format words or short inline content.
    \item The browser displays the content written inside \texttt{<body>} and \texttt{</body>}.
\end{itemize}

\begin{summarybox}
The \texttt{<head>} holds information about the page, while the \texttt{<body>} holds what users see.
\end{summarybox}

% ─────────────────────────────────────────
\section{HTML Tags}
% ─────────────────────────────────────────

\begin{definitionbox}[Tag]
A tag is a keyword enclosed in angle brackets used to define elements in HTML.
\end{definitionbox}

\subsection{Types of Tags}

\begin{itemize}
    \item \term{Container (Paired) Tags}: Have opening and closing parts.  
          Example: \texttt{<p> ... </p>}
    \item \term{Empty (Unpaired) Tags}: Do not have a closing tag.  
          Example: \texttt{<br>}, \texttt{<hr>}
\end{itemize}

In modern HTML writing, empty elements may also be written with a closing slash, such as \texttt{<br/>} or \texttt{<hr/>}. Closing paired tags properly is important because missing end tags can produce unexpected output in a browser.

\subsection{Examples}

This example compares paired tags and empty tags.
A paragraph uses both opening and closing tags, but \texttt{br} and \texttt{hr} are empty elements.

\begin{lstlisting}[style=htmlide, caption={Paired and Empty HTML Tags}]
<p>This is a paragraph</p>
<br>
<hr>
\end{lstlisting}

\textbf{What to notice}
\begin{itemize}
    \item \texttt{p} is a container tag because it wraps content.
    \item \texttt{br} creates a line break without a closing tag.
    \item \texttt{hr} creates a horizontal line.
\end{itemize}

\begin{summarybox}
Tags define how content is structured and displayed in a webpage.
\end{summarybox}

\subsection{HTML Comments}

Comments are used to make HTML code easier to understand. Browsers ignore comments, so they are not displayed on the webpage.

\begin{lstlisting}[style=htmlide, caption={HTML Comment Syntax}]
<!-- This is a comment -->
\end{lstlisting}

\textbf{What to notice}
\begin{itemize}
    \item Comments are useful for explaining code.
    \item Comments are visible in source code but not on the rendered webpage.
\end{itemize}

% ─────────────────────────────────────────
\section{HTML Elements}
% ─────────────────────────────────────────

\begin{definitionbox}[Element]
An HTML element consists of a start tag, content, and an end tag.
\end{definitionbox}

\subsection{Structure of an Element}

\begin{formulabox}
\[
\text{Element} = \text{Start Tag} + \text{Content} + \text{End Tag}
\]
\end{formulabox}

\subsection{Example}

This short example shows one complete HTML element.
The start tag, content, and end tag together form the element.

\begin{lstlisting}[style=htmlide, caption={A Complete Paragraph Element}]
<p>Hello World</p>
\end{lstlisting}

\textbf{What to notice}
\begin{itemize}
    \item The content is written between the two tags.
    \item The closing tag uses a slash before the tag name.
\end{itemize}

\begin{itemize}
    \item \texttt{<p>}: Start tag
    \item \texttt{Hello World}: Content
    \item \texttt{</p>}: End tag
\end{itemize}

\begin{summarybox}
Tags are the markers, while elements are the complete units formed by tags and content.
\end{summarybox}

% ─────────────────────────────────────────
\section{HTML Attributes}
% ─────────────────────────────────────────

\begin{definitionbox}[Attribute]
An attribute provides additional information about an HTML element and is written inside the start tag.
\end{definitionbox}

\subsection{General Form}

\begin{formulabox}
\[
\text{<tag attribute="value">}
\]
\end{formulabox}

\subsection{Examples}

This example shows how attributes add extra information to HTML elements.
The \texttt{href} attribute gives a link destination, while \texttt{src} and \texttt{alt} describe an image source and fallback text.

\begin{lstlisting}[style=htmlide, caption={Common HTML Attributes}]
<a href="https://www.example.com">Visit</a>
<img src="image.jpg" alt="Image">
\end{lstlisting}

\textbf{What to notice}
\begin{itemize}
    \item Attributes are written inside the opening tag.
    \item Attribute values are usually written inside quotation marks.
    \item The same element can have more than one attribute.
\end{itemize}

\subsection{Common Attributes}

\begin{itemize}
    \item \texttt{href}: defines link destination
    \item \texttt{src}: defines image source
    \item \texttt{alt}: alternative text for images
    \item \texttt{title}: tooltip text
\end{itemize}

\begin{summarybox}
Attributes modify the behavior or appearance of elements.
\end{summarybox}

% ─────────────────────────────────────────
\section{Headings, Paragraphs and Line Breaks}

\subsection{Heading Elements}

HTML provides six levels of headings from \texttt{h1} to \texttt{h6}. They are container elements and are used to divide text into sections.

\begin{itemize}
    \item \texttt{h1} is the largest and most important heading.
    \item \texttt{h6} is the smallest heading.
    \item Headings help readers and search engines understand document structure.
\end{itemize}

\begin{lstlisting}[style=htmlide, caption={Heading Elements}]
<h1>Main Heading</h1>
<h2>Sub Heading</h2>
<h3>Smaller Heading</h3>
\end{lstlisting}

\textbf{What to notice}
\begin{itemize}
    \item Heading tags should be used for structure, not only for large text.
    \item Each heading has an opening tag and a closing tag.
\end{itemize}

\subsection{Paragraph, Break and Horizontal Rule}

The \texttt{p} tag defines a paragraph. Browsers collapse extra spaces and line breaks in source code unless specific HTML tags are used.

\begin{itemize}
    \item \texttt{<p>...</p>} creates a paragraph.
    \item \texttt{<br/>} creates a line break.
    \item \texttt{<hr/>} creates a horizontal line.
\end{itemize}

\begin{lstlisting}[style=htmlide, caption={Paragraph and Line Break}]
<p>This is a paragraph.</p>
<p>This line has<br/>a manual break.</p>
<hr/>
\end{lstlisting}

\textbf{What to notice}
\begin{itemize}
    \item Pressing Enter in source code does not always create a new browser line.
    \item Use \texttt{br} for line breaks and \texttt{p} for paragraphs.
\end{itemize}

\section{Text Formatting Tags}
% ─────────────────────────────────────────

\subsection{Basic Formatting}

\begin{longtable}{|C{4cm}|C{6cm}|C{5cm}|}
\longtableheader{3}{Formatting Tags}{%
\tblcolhdcell{Tag} & \tblcolhdcell{Purpose} & \tblcolhdcell{Example}}

<b> & Bold text & \texttt{<b>Text</b>} \\ \tblhline
<i> & Italic text & \texttt{<i>Text</i>} \\ \tblhline
<u> & Underline text & \texttt{<u>Text</u>} \\ \tblhline
<sup> & Superscript & \texttt{a<sup>2</sup>} \\ \tblhline
<sub> & Subscript & \texttt{H<sub>2</sub>O} \\ \tblhline
<br> & Line break & \texttt{Line<br>Next} \\ \tblhline
<hr> & Horizontal line & \texttt{<hr>} \\ \tblhline
<small> & Smaller text & \texttt{<small>Text</small>} \\ \tblhline
<strong> & Important bold text & \texttt{<strong>Text</strong>} \\ \tblhline
<em> & Emphasized italic text & \texttt{<em>Text</em>} \\ \tblhline
<mark> & Highlighted text & \texttt{<mark>Text</mark>} \\ \tblhline
<del> & Deleted text & \texttt{<del>Text</del>} \\ \tblhline
<ins> & Inserted text & \texttt{<ins>Text</ins>} \\ \tblhline

\tblbottomrule
\caption{Common HTML Formatting Tags}
\end{longtable}

\subsection{Font, Bangla Text and Marquee}

\begin{itemize}
    \item The old \texttt{font} tag was used to set font face, color and size. It is exam/reference only, not modern practice.
    \item Modern webpages should use CSS, such as the \texttt{style} attribute or stylesheet, for font family, size and color.
    \item For Bangla text, Unicode with UTF-8 encoding is the correct modern method.
    \item Old font names such as \texttt{SutonnyMJ} may appear in older examples, but they are not recommended for modern web pages.
    \item The \texttt{marquee} tag creates scrolling text. It is exam/reference only, not modern practice.
\end{itemize}

\begin{lstlisting}[style=htmlide, caption={Modern Font Style with Unicode Text}]
<p style="font-family: Arial; color: red;">
    This is styled text.
</p>
<p>Bangla Unicode text</p>
\end{lstlisting}

\textbf{What to notice}
\begin{itemize}
    \item CSS is preferred instead of the old \texttt{font} tag.
    \item UTF-8 and Unicode make Bangla text portable across modern browsers.
\end{itemize}

\begin{summarybox}
Formatting tags are used to control the appearance of text.
\end{summarybox}

% ─────────────────────────────────────────
\section{HTML Lists}
% ─────────────────────────────────────────

HTML provides different types of lists to organize content.

\subsection{Ordered List (OL)}

\begin{definitionbox}[Ordered List]
An ordered list displays items in a numbered sequence.
\end{definitionbox}

This example creates a numbered list.
Each list item must be written inside an \texttt{li} element.

\begin{lstlisting}[style=htmlide, caption={Ordered List}]
<ol>
    <li>Item 1</li>
    <li>Item 2</li>
</ol>
\end{lstlisting}

\textbf{What to notice}
\begin{itemize}
    \item \texttt{ol} means ordered list.
    \item Browsers number ordered list items automatically.
\end{itemize}

\subsection{Unordered List (UL)}

\begin{definitionbox}[Unordered List]
An unordered list displays items using bullet points.
\end{definitionbox}

This example creates a bullet list.
It is useful when the order of items is not important.

\begin{lstlisting}[style=htmlide, caption={Unordered List}]
<ul>
    <li>Item A</li>
    <li>Item B</li>
</ul>
\end{lstlisting}

\textbf{What to notice}
\begin{itemize}
    \item \texttt{ul} means unordered list.
    \item Each bullet is created by an \texttt{li} tag.
\end{itemize}

\subsection{Description List (DL)}

\begin{definitionbox}[Description List]
A description list is used to define terms and their descriptions.
\end{definitionbox}

This example creates a description list.
It pairs a term with a short explanation.

\begin{lstlisting}[style=htmlide, caption={Description List}]
<dl>
    <dt>HTML</dt>
    <dd>Markup language</dd>
</dl>
\end{lstlisting}

\textbf{What to notice}
\begin{itemize}
    \item \texttt{dt} defines the term.
    \item \texttt{dd} gives the description of that term.
\end{itemize}

\begin{summarybox}
HTML lists help organize content in a structured and readable way.
\end{summarybox}

% ─────────────────────────────────────────
\section{Hyperlinks in HTML}
% ─────────────────────────────────────────

\begin{definitionbox}[Hyperlink]
A hyperlink is a connection that allows navigation from one webpage to another.
\end{definitionbox}

\subsection{Anchor Tag}

This example creates a clickable link.
The \texttt{href} attribute stores the destination address, and the text between the tags becomes clickable.

\begin{lstlisting}[style=htmlide, caption={Basic Hyperlink}]
<a href="https://www.example.com">Visit Website</a>
\end{lstlisting}

\textbf{What to notice}
\begin{itemize}
    \item The \texttt{a} tag creates an anchor.
    \item The link destination is controlled by \texttt{href}.
\end{itemize}

\subsection{Important Attributes}

\begin{itemize}
    \item \texttt{href}: destination URL
    \item \texttt{target}: where to open link
\end{itemize}

\subsection{Target Values}

\begin{itemize}
    \item \texttt{\_self}: opens in same tab
    \item \texttt{\_blank}: opens in new tab
\end{itemize}

\begin{summarybox}
Hyperlinks are the foundation of web navigation, connecting pages together.
\end{summarybox}

% ─────────────────────────────────────────
\section{Images in HTML}
% ─────────────────────────────────────────

\begin{definitionbox}[Image Tag]
The \texttt{<img>} tag is used to display images on a webpage.
\end{definitionbox}

\subsection{Basic Syntax}

This example inserts an image into a webpage.
The \texttt{src} attribute gives the image file path, and \texttt{alt} provides text if the image cannot load.

\begin{lstlisting}[style=htmlide, caption={Basic Image Tag}]
<img src="image.jpg" alt="Description">
\end{lstlisting}

\textbf{What to notice}
\begin{itemize}
    \item \texttt{img} is an empty element.
    \item \texttt{alt} improves accessibility and helps when the image path is wrong.
\end{itemize}

\subsection{Common Attributes}

\begin{itemize}
    \item \texttt{src}: specifies image path
    \item \texttt{alt}: alternative text if image fails to load
    \item \texttt{width}: image width
    \item \texttt{height}: image height
\end{itemize}

\begin{summarybox}
Images enhance visual presentation and improve user experience.
\end{summarybox}

% ─────────────────────────────────────────
\section{HTML Tables}
% ─────────────────────────────────────────

\begin{definitionbox}[Table]
An HTML table is used to display data in rows and columns.
\end{definitionbox}

\subsection{Basic Structure}

This example builds a small table using rows, header cells, and data cells.
The first row contains headings, and the second row contains normal table data.

\begin{lstlisting}[style=htmlide, caption={Basic HTML Table}]
<table border="1">
    <tr>
        <th>Name</th>
        <th>Age</th>
    </tr>
    <tr>
        <td>John</td>
        <td>20</td>
    </tr>
</table>
\end{lstlisting}

\textbf{What to notice}
\begin{itemize}
    \item \texttt{tr} creates a table row.
    \item \texttt{th} creates a heading cell.
    \item \texttt{td} creates a normal data cell.
\end{itemize}

\subsection{Table Elements}

\begin{itemize}
    \item \texttt{<table>}: defines table
    \item \texttt{<tr>}: table row
    \item \texttt{<th>}: table header cell
    \item \texttt{<td>}: table data cell
\end{itemize}

\subsection{Common Attributes}

\begin{itemize}
    \item \texttt{border}: defines border size
    \item \texttt{cellpadding}: space inside cells
    \item \texttt{cellspacing}: space between cells
\end{itemize}

\begin{summarybox}
Tables are used to present structured data clearly in rows and columns.
\end{summarybox}


% ─────────────────────────────────────────
\section{Final Summary of the Chapter}
% ─────────────────────────────────────────

\begin{summarybox}[Chapter Overview]
\begin{itemize}
    \item Web design deals with creating and structuring websites.
    \item Websites are collections of interconnected webpages.
    \item IP address uniquely identifies devices on a network.
    \item Domain name provides a human-readable address.
    \item URL specifies the location of resources.
    \item HTTP controls browser-server webpage communication.
    \item FTP is used for uploading and downloading files.
    \item Website structure affects navigation and usability.
    \item Websites may be static, dynamic, local or remote.
    \item Websites can also be classified by purpose, such as business, archive, news or e-commerce sites.
    \item Browsers interpret webpage code, while search engines help find webpages.
    \item HTML is used to create webpage structure.
    \item HTML documents contain head and body sections.
    \item Tags, elements, and attributes form HTML content.
    \item Lists, images, links, and tables are essential HTML components.
\end{itemize}
\end{summarybox}